\scanpage{91}
% loc 1–2
\[
u'=(\operatorname{Tr}u)1-u=b.1-u
\qquad\left\{\begin{aligned}\operatorname{Tr}u'&=\operatorname{Tr}u\\ \det u'&=\deleted{\unsure{\operatorname{Tr}v}}\det u\end{aligned}\right.
\]
% loc 3–5
alors
\[\begin{aligned}
\tau'&=\operatorname{Tr}u'v=\operatorname{Tr}\deleted{\unsure{u}}(b-u)v=b\operatorname{Tr}v)-\operatorname{Tr}(uv)\\
&=bb'-\tau\deleted{\unsure{(uv)}}
\end{aligned}\]
donc \deleted{\unsure{$t'$}} $\tau'=t$, On gagne !

% loc 6
Que se passe-t-il si
% loc 7
\[ (b^2-4c)(b'^2-4c')\in k^{\unsure{\times|*}} \tag*{$(12_1)$} \]
% loc 8
\deleted{mais} \inserted{Si} $M(b,c,b',c',t)$ est azumayen, je dis que
% loc 9
\[ \underbrace{t^2-bb't+b^2c'+b'^2c-4cc'}_{Q(t;b,c,b',c')\deleted{\unsure{t}}}\in k^*\quad ? \tag*{$(12_2)$} \]
% loc 10–11
\deleted{Je dis qu'alors $M(b,c,b',c',t)$ \uline{n'est pas} \unsure{Azumayen|Azumaya|Azumayenne}.} En effet, \deleted{si c'était}
% loc 12–17
Comme $A_1=k[\underline u]$ et $A_2=k[\underline v]$ \deleted{\unsure{q}} sont étales, et en position générale, il \deleted{les formules} y aurait par la théorie générale un invariant $\rho_0$, relié à $b,c,b',c',t$ par $(11')$, récrit sous forme
% loc 18
\[ \deleted{\rho_0(t^2-bb't}\qquad \rho_0Q=\delta\delta' \]
% loc 19
Comme $\delta\delta'\in k^*$, on doit avoir $Q\in k^*$

% loc 20–21
\uline{Remarque} Le discriminant de $Q$ (en tant que polynôme en $t$) est
% archive sheet number: 90
